\chapter{Appendix: Matrices and Discrete Math}
\label{ch39:matrices}

This appendix is a toolbox. Throughout the book we leaned on a handful
of mathematical objects -- matrices, graphs, logarithms, factorials -- as if you
already had them in your pocket. Here we slow down and build each one from the
ground up, just thoroughly enough that nothing in the main chapters has to feel
like magic. Read it front to back, or keep it on the shelf and
flip to the piece you need.

None of it is hard. A matrix is a grid of numbers with
some bookkeeping rules. A graph is dots and lines. A logarithm answers
``how many times did I halve?'' For each idea we attach a
picture, work a number example by hand, and then state the thing precisely so you can
trust it later.

\begin{focusbox}
This appendix maps to CLRS Appendices C and D (and a little of A and B). When
CLRS writes ``recall from Appendix D that matrix multiplication is associative,''
this is the chapter that earns you the right to recall it.
\end{focusbox}

\begin{historybox}
For centuries, mathematicians solved linear systems of equations using ad-hoc
substitutions or cumbersome determinant formulas, viewing numbers as individual
scalars. The leap to treating a whole rectangular grid of numbers as a single
algebraic object came in 1858 from the English mathematician Arthur Cayley in his
landmark paper \emph{A Memoir on the Theory of Matrices}. Cayley defined matrix
addition, scalar multiplication, matrix multiplication, and inverses, proving that
matrices form a closed, associative algebraic system. Cayley's formulation
liberated linear equations from scalar drudgery, laying the groundwork for modern
linear algebra.

Half a century earlier, the computational power of matrix reduction had already
rescued astronomy. On New Year's Day in 1801, Italian astronomer Giuseppe Piazzi
discovered the dwarf planet Ceres, tracking it for only forty days before it was
lost in the daytime glare of the Sun. Astronomers across Europe failed to predict
where it would re-emerge from behind the Sun. A twenty-four-year-old Carl Friedrich
Gauss took up the challenge. Developing the method of least squares and systematic
row elimination---what we now call \emph{Gaussian elimination}---Gauss solved the
massive system of linear equations governing orbital parameters. When Ceres
reappeared in December 1801, it was precisely where Gauss had predicted. In his
1809 treatise \emph{Theoria Motus}, Gauss codified the elimination algorithm that
remains the fundamental numerical workhorse of scientific computing.

Underneath Cayley's algebra and Gauss's elimination sits a profound geometric
truth: the \emph{determinant}. Rather than a dry algebraic formula, the
determinant $\det(A)$ measures the factor by which the linear transformation $A$
scales hypervolumes in space. In two dimensions, $\det(A)$ is the factor by which
$A$ stretches or squashes area; its sign indicates whether the transformation
preserves or reverses spatial orientation. A determinant of zero means space has
been flattened into a lower dimension---a line or a point---irreversibly destroying
information and explaining why matrices with zero determinant cannot be inverted.
\end{historybox}

\section{Matrices: grids of numbers that \emph{do} something}
\label{ch39:matrices-intro}

\subsection{What a matrix is}

A \emph{matrix} is a rectangular grid of numbers. We describe its shape by saying
``$m$ by $n$,'' written $m \times n$, meaning $m$ rows and $n$ columns. Here is a
$2 \times 3$ matrix (two rows, three columns):
\[
A \;=\;
\begin{pmatrix}
2 & 0 & -1 \\
5 & 3 & \phantom{-}4
\end{pmatrix}.
\]
The number in row $i$, column $j$ is called the $(i,j)$ \emph{entry} and written
$A_{ij}$ (or $a_{ij}$). So $A_{11} = 2$, $A_{13} = -1$, and $A_{23} = 4$. Always
read the row index first, then the column. That ordering is a convention you must
keep straight, because $A_{23}$ and $A_{32}$ are different requests entirely (and
in our example $A_{32}$ does not even exist -- there is no third row).

\begin{intuitionbox}
Picture a spreadsheet. A matrix is exactly a spreadsheet of numbers with the
column headers and row labels stripped away, leaving just the cells. Everything
we do to matrices is something you could imagine doing to a spreadsheet: add two
sheets cell by cell, scale every cell, or combine two sheets with a more
elaborate row-meets-column rule. Hold onto the spreadsheet picture and you will
rarely get lost.
\end{intuitionbox}

A matrix with the same number of rows and columns ($m = n$) is called
\emph{square}. A matrix with a single column ($n = 1$) is a \emph{column vector};
one with a single row ($m = 1$) is a \emph{row vector}. We will write vectors as
lowercase letters, like $x$, and matrices as uppercase, like $A$.

\begin{focusbox}
Three numbers fully pin down the ``shape situation'' of any matrix product, and
keeping them visible saves you from most mistakes: the matrix is $m \times n$,
the row index $i$ runs $1 \le i \le m$, and the column index $j$ runs
$1 \le j \le n$. Write the shape next to every matrix while you are learning.
\end{focusbox}

\subsection{Adding matrices and scaling them}

Two matrices of the \emph{same shape} are added entry by entry. If $A$ and $B$
are both $m \times n$, then $A + B$ is the $m \times n$ matrix whose $(i,j)$ entry
is $A_{ij} + B_{ij}$. There is no cleverness here at all:
\[
\begin{pmatrix} 2 & 0 \\ 5 & 3 \end{pmatrix}
+
\begin{pmatrix} 1 & 4 \\ -2 & 1 \end{pmatrix}
=
\begin{pmatrix} 2{+}1 & 0{+}4 \\ 5{-}2 & 3{+}1 \end{pmatrix}
=
\begin{pmatrix} 3 & 4 \\ 3 & 4 \end{pmatrix}.
\]
If the shapes do not match, the sum is undefined -- there is no row to pair
with a missing row. \emph{Scaling} (multiplying by a single number $c$, called a
\emph{scalar}) multiplies every entry by $c$:
\[
3 \cdot \begin{pmatrix} 2 & 0 \\ 5 & 3 \end{pmatrix}
= \begin{pmatrix} 6 & 0 \\ 15 & 9 \end{pmatrix}.
\]
Addition is commutative ($A + B = B + A$) and associative
($(A+B)+C = A+(B+C)$) for the same reason ordinary number addition is: it all
happens entry by entry, and the entries are just numbers.

\subsection{The product: where matrices start pulling their weight}
\label{ch39:matmul}

Matrix \emph{multiplication} is the one rule that looks strange the first time and
obvious the tenth time. To multiply $A$ by $B$ and form $C = AB$, the shapes must
\emph{interlock}: the number of columns of $A$ must equal the number of rows of
$B$. If $A$ is $m \times n$ and $B$ is $n \times p$, then $C = AB$ is
$m \times p$, and
\[
C_{ij} \;=\; \sum_{k=1}^{n} A_{ik}\,B_{kj}.
\]
In words: the $(i,j)$ entry of the product is the \emph{dot product} of row $i$
of $A$ with column $j$ of $B$. You march across the row and down the column in
lockstep, multiply the paired numbers, and add up the results.

\begin{intuitionbox}
The shared dimension $n$ is the ``handshake.'' Row $i$ of $A$ reaches out with $n$
fingers; column $j$ of $B$ reaches back with $n$ fingers; each finger meets its
partner, they multiply, and you total the handshake. If the finger counts do not
match, no handshake is possible -- that is exactly why the inner dimensions must
agree.
\end{intuitionbox}

Let us do one slowly. Take
\[
A = \begin{pmatrix} 1 & 2 \\ 0 & 3 \end{pmatrix}
\quad (2 \times 2),
\qquad
B = \begin{pmatrix} 4 & 5 \\ 6 & 7 \end{pmatrix}
\quad (2 \times 2).
\]
The inner dimensions match ($2 = 2$), so $C = AB$ is $2 \times 2$. Compute each
entry as a row-times-column handshake:
\begin{align*}
C_{11} &= (1)(4) + (2)(6) = 4 + 12 = 16, \\
C_{12} &= (1)(5) + (2)(7) = 5 + 14 = 19, \\
C_{21} &= (0)(4) + (3)(6) = 0 + 18 = 18, \\
C_{22} &= (0)(5) + (3)(7) = 0 + 21 = 21.
\end{align*}
So
\[
AB = \begin{pmatrix} 16 & 19 \\ 18 & 21 \end{pmatrix}.
\]

\begin{labbox}
Trace it yourself. Compute $BA$ for the same two matrices and compare it with
$AB$ above. You should get
$BA = \begin{psmallmatrix} 4 & 23 \\ 6 & 33 \end{psmallmatrix}$,
which is \emph{not} equal to $AB$. Sit with that for a second: the order in which
you multiply matrices matters. We will turn this observation into a rule below.
\end{labbox}

\begin{trapbox}
The single most common matrix mistake is assuming $AB = BA$. \textbf{Matrix
multiplication is not commutative.} It is also easy to forget the shape rule and
try to multiply, say, a $2 \times 3$ by a $2 \times 2$: the inner dimensions
($3$ and $2$) disagree, so the product does not exist. Always check the handshake
before you compute a single entry.
\end{trapbox}

\subsection{A picture of the handshake}

The diagram below shows how the highlighted row of $A$ and the highlighted column
of $B$ produce a single entry of $C$.

\begin{center}
\begin{tikzpicture}[scale=0.9, every node/.style={font=\small}]
  % Matrix A
  \node at (0,2.4) {$A$};
  \draw (-0.7,0.7) rectangle (1.3,2.0);
  \draw (-0.7,1.35) -- (1.3,1.35);
  \draw (0.3,0.7) -- (0.3,2.0);
  \node[sourceblue] at (-0.2,1.675) {$1$};
  \node[sourceblue] at (0.8,1.675) {$2$};
  \node at (-0.2,1.025) {$0$};
  \node at (0.8,1.025) {$3$};
  \draw[sourceblue,thick,rounded corners] (-0.65,1.4) rectangle (1.25,1.95);

  % Matrix B
  \node at (3.3,2.4) {$B$};
  \draw (2.6,0.7) rectangle (4.6,2.0);
  \draw (2.6,1.35) -- (4.6,1.35);
  \draw (3.6,0.7) -- (3.6,2.0);
  \node[targetgreen] at (3.1,1.675) {$4$};
  \node at (4.1,1.675) {$5$};
  \node[targetgreen] at (3.1,1.025) {$6$};
  \node at (4.1,1.025) {$7$};
  \draw[targetgreen,thick,rounded corners] (2.65,0.75) rectangle (3.55,1.95);

  % Arrow to result
  \draw[-{Latex},thick] (4.9,1.35) -- (5.8,1.35);

  % Result C entry
  \node at (7.3,2.4) {$C = AB$};
  \draw (6.3,0.7) rectangle (8.3,2.0);
  \draw (6.3,1.35) -- (8.3,1.35);
  \draw (7.3,0.7) -- (7.3,2.0);
  \node[structurepurple,font=\bfseries] at (6.8,1.675) {$16$};
  \node at (7.8,1.675) {$19$};
  \node at (6.8,1.025) {$18$};
  \node at (7.8,1.025) {$21$};
  \draw[structurepurple,thick,rounded corners] (6.35,1.4) rectangle (7.25,1.95);

  \node[align=center,font=\footnotesize] at (3.6,0.1)
    {$C_{11} = (1)(4) + (2)(6) = 16$};
\end{tikzpicture}
\end{center}

The blue row of $A$ shakes hands with the green column of $B$: $(1)(4) + (2)(6) =
16$, which lands in the purple top-left cell of $C$. Every other entry is the same
ritual with a different row and column.

\subsection{The identity matrix and the transpose}

The \emph{identity matrix} $I_n$ is the $n \times n$ square matrix with $1$'s down
the main diagonal (the entries $I_{ii}$) and $0$'s everywhere else. For example,
\[
I_3 = \begin{pmatrix} 1 & 0 & 0 \\ 0 & 1 & 0 \\ 0 & 0 & 1 \end{pmatrix}.
\]
It earns its name: for any $m \times n$ matrix $A$ we have $I_m A = A$ and
$A I_n = A$. Multiplying by $I$ changes nothing, exactly as multiplying a number
by $1$ changes nothing. The identity is the ``do-nothing'' matrix, and it is the
multiplicative anchor for everything else.

The \emph{transpose} of $A$, written $A^{\mathsf T}$, flips the matrix across its
main diagonal: row $i$ becomes column $i$. Formally $(A^{\mathsf T})_{ij} =
A_{ji}$. If $A$ is $m \times n$, then $A^{\mathsf T}$ is $n \times m$:
\[
A = \begin{pmatrix} 2 & 0 & -1 \\ 5 & 3 & 4 \end{pmatrix}
\;\Longrightarrow\;
A^{\mathsf T} = \begin{pmatrix} 2 & 5 \\ 0 & 3 \\ -1 & 4 \end{pmatrix}.
\]
A handy fact you will meet again: the transpose of a product reverses the order,
$(AB)^{\mathsf T} = B^{\mathsf T} A^{\mathsf T}$. A matrix equal to its own
transpose ($A = A^{\mathsf T}$, which forces it to be square) is called
\emph{symmetric}; these show up constantly, for instance as adjacency matrices of
undirected graphs later in this chapter.

\subsection{The algebra rules}

Here is the short list of laws you may assume in the rest of the book. They are
not difficult to prove -- each follows from pushing the sum
$\sum_k A_{ik} B_{kj}$ around -- but they are worth seeing collected in one place.

\begin{formalbox}
For matrices of compatible shapes and any scalar $c$:
\begin{itemize}
  \item \textbf{Associative:} $(AB)C = A(BC)$.
  \item \textbf{Distributive:} $A(B+C) = AB + AC$ and $(A+B)C = AC + BC$.
  \item \textbf{Scalar slide:} $c(AB) = (cA)B = A(cB)$.
  \item \textbf{Identity:} $IA = A = AI$ (for square $A$).
  \item \textbf{Transpose of a product:} $(AB)^{\mathsf T} = B^{\mathsf T}A^{\mathsf T}$.
  \item \textbf{NOT commutative:} in general $AB \ne BA$.
\end{itemize}
\end{formalbox}

The associative law deserves a moment, because it is doing quiet work in several
chapters. It says that when you multiply a chain $A_1 A_2 \cdots A_r$, the
\emph{answer} does not depend on how you parenthesize the multiplications -- only
the \emph{cost} of computing it does. That gap between ``same answer, different
cost'' is the entire reason the matrix-chain multiplication problem (a classic
dynamic-programming example) is interesting.

\begin{aestheticbox}
Why carry matrices into an algorithms book at all? Because they are the language
of \emph{everything done in bulk}. A single matrix multiply applies one linear
transformation to a whole batch of vectors at once; that is what powers 3D
graphics pipelines, the layers of a neural network, Google's original PageRank
computation, and the fast all-pairs-shortest-paths trick that reuses matrix
multiplication over a funny ``min-plus'' arithmetic. When you can phrase a
computation as a matrix product, decades of fast multiplication algorithms (and
fast hardware) become available to you for free.
\end{aestheticbox}

\subsection{The cost of multiplying matrices}
\label{ch39:matmul-cost}

How expensive is one matrix multiply? Take the square case, $A$ and $B$ both
$n \times n$, so $C = AB$ is also $n \times n$. The product has $n^2$ entries.
Each entry $C_{ij}$ is a handshake of length $n$: $n$ multiplications and $n - 1$
additions. So the total arithmetic is
\[
\underbrace{n^2}_{\text{entries}} \times \underbrace{(2n - 1)}_{\text{ops each}}
\;=\; 2n^3 - n^2 \;=\; \Theta(n^3).
\]
The naive algorithm is three nested loops, and counting them gives the same
answer directly.

\begin{algorithm}[H]
\caption{Naive square matrix multiplication}
\KwIn{$n \times n$ matrices $A$ and $B$}
\KwOut{$C = AB$}
\For{$i \leftarrow 1$ \KwTo $n$}{
  \For{$j \leftarrow 1$ \KwTo $n$}{
    $s \leftarrow 0$\;
    \For{$k \leftarrow 1$ \KwTo $n$}{
      $s \leftarrow s + A_{ik}\cdot B_{kj}$\;
    }
    $C_{ij} \leftarrow s$\;
  }
}
\Return $C$\;
\end{algorithm}

The body $s \leftarrow s + A_{ik} B_{kj}$ runs exactly $n \cdot n \cdot n = n^3$
times, once for each triple $(i,j,k)$, so the running time is $\Theta(n^3)$.

\begin{complexitybox}
For two $n \times n$ matrices, the naive algorithm above costs:
\begin{itemize}
  \item Time: $\Theta(n^3)$ arithmetic operations.
  \item Extra space: $\Theta(n^2)$ to store the output (or $\Theta(1)$ extra
        beyond the output, since each $C_{ij}$ is computed independently).
\end{itemize}
Strassen's algorithm (CLRS Chapter 4) beats this with $\Theta(n^{\log_2 7})
\approx \Theta(n^{2.81})$ by trading some multiplications for additions via
divide and conquer. The current theoretical record is below $n^{2.372}$, though
nobody uses those galactic algorithms in practice.
\end{complexitybox}

\begin{walkbox}
On your walk, picture the three nested loops as a $n \times n \times n$ cube of
little ``multiply-and-add'' events. The naive algorithm visits every one of the
$n^3$ cells. Strassen's insight is that those cells are not all independent --
clever combinations let you skip about an eighth of the work at each level of
recursion. Can you feel, intuitively, why \emph{recursion} (splitting each matrix
into four quadrants) is the natural place to look for such savings?
\end{walkbox}

\section{A pinch of linear algebra}
\label{ch39:linalg}

You do not need a full linear-algebra course to read this book, but a few facts
recur. We collect them here with the lightest possible touch -- enough to
recognize the words, not to pass a final exam.

\subsection{Vectors, linear combinations, and independence}

A \emph{vector} in this appendix is just a column matrix of $n$ numbers, an
arrow's worth of data. Given vectors $v_1, \dots, v_k$ and scalars
$c_1, \dots, c_k$, the expression
\[
c_1 v_1 + c_2 v_2 + \cdots + c_k v_k
\]
is called a \emph{linear combination} of them. It is the most natural way to
``mix'' vectors: scale each one, then add. The set of \emph{all} linear
combinations of a fixed collection of vectors is called their \emph{span}.

The vectors are \emph{linearly independent} if the only way to make the zero
vector as a linear combination is the boring way, with every coefficient zero. If
some non-trivial combination collapses to zero, the vectors are
\emph{dependent} -- one of them is redundant, expressible from the others. For
example, $(1,0)$, $(0,1)$, and $(2,3)$ are dependent because
$2(1,0) + 3(0,1) - (2,3) = (0,0)$ with not-all-zero coefficients.

\begin{intuitionbox}
Independence means ``no freeloaders.'' Each vector contributes a genuinely new
direction that the others cannot reproduce by mixing. Dependence means at least
one vector is a freeloader -- you could delete it and lose nothing, because the
remaining vectors can already reach it.
\end{intuitionbox}

\subsection{Determinant and inverse, briefly}

For a square matrix $A$, the \emph{determinant} $\det(A)$ is a single number
squeezed out of the whole grid that tells you, among other things, whether $A$ is
\emph{invertible}. For the $2 \times 2$ case the formula is short and worth
memorizing:
\[
\det\!\begin{pmatrix} a & b \\ c & d \end{pmatrix} = ad - bc.
\]
A square matrix $A$ has an \emph{inverse} $A^{-1}$ -- a matrix with
$A A^{-1} = A^{-1} A = I$ -- exactly when $\det(A) \ne 0$. The inverse is the
matrix analogue of a reciprocal: $A^{-1}$ undoes whatever $A$ does. For the
$2 \times 2$ case, when $ad - bc \ne 0$,
\[
\begin{pmatrix} a & b \\ c & d \end{pmatrix}^{-1}
= \frac{1}{ad - bc}\begin{pmatrix} d & -b \\ -c & a \end{pmatrix}.
\]
Let us check one. With $A = \begin{psmallmatrix} 2 & 1 \\ 1 & 1 \end{psmallmatrix}$
we have $\det(A) = 2\cdot 1 - 1 \cdot 1 = 1$, so
$A^{-1} = \begin{psmallmatrix} 1 & -1 \\ -1 & 2 \end{psmallmatrix}$. Multiply to
confirm:
\[
\begin{pmatrix} 2 & 1 \\ 1 & 1 \end{pmatrix}
\begin{pmatrix} 1 & -1 \\ -1 & 2 \end{pmatrix}
= \begin{pmatrix} 2{-}1 & -2{+}2 \\ 1{-}1 & -1{+}2 \end{pmatrix}
= \begin{pmatrix} 1 & 0 \\ 0 & 1 \end{pmatrix} = I. \checkmark
\]

\begin{trapbox}
A square matrix with $\det(A) = 0$ is called \emph{singular} and has \emph{no}
inverse at all -- you cannot ``divide'' by it. This is the matrix world's version
of dividing by zero, and the determinant is precisely the warning light. Whenever
an algorithm hands you a matrix it intends to invert (solving linear systems,
least-squares fitting), a zero determinant is the place where things break.
\end{trapbox}

\begin{figure}[H]
\centering
\begin{tikzpicture}[
  scale=0.88,
  >=Stealth,
  lbl/.style={font=\sffamily\footnotesize}
]
% Panel (a): Non-singular transformation (det > 0)
\begin{scope}[shift={(0,0)}]
  % Axes for input unit square
  \draw[->, gray!60] (-0.3,0) -- (2.0,0) node[right, lbl, black]{\scriptsize $x$};
  \draw[->, gray!60] (0,-0.3) -- (0,2.0) node[above, lbl, black]{\scriptsize $y$};

  \fill[sourceblue!20] (0,0) rectangle (1.2,1.2);
  \draw[->, thick, sourceblue!80!black] (0,0) -- (1.2,0) node[below, lbl]{$\hat{i}$};
  \draw[->, thick, sourceblue!80!black] (0,0) -- (0,1.2) node[left, lbl]{$\hat{j}$};
  \node[lbl] at (0.6,0.6) {Area $= 1$};

  % Transform arrow
  \draw[->, very thick, quotientorange] (2.2, 1.0) -- (3.2, 1.0)
    node[midway, above, lbl, font=\scriptsize\bfseries]{$A = \left[\begin{smallmatrix} 2 & 1 \\ 1 & 2 \end{smallmatrix}\right]$};

  % Output area
  \begin{scope}[shift={(3.8,0)}]
    \draw[->, gray!60] (-0.3,0) -- (3.4,0) node[right, lbl, black]{\scriptsize $x'$};
    \draw[->, gray!60] (0,-0.3) -- (0,3.2) node[above, lbl, black]{\scriptsize $y'$};

    \coordinate (Ai) at (2.0, 1.0);
    \coordinate (Aj) at (1.0, 2.0);
    \coordinate (Aij) at (3.0, 3.0);

    \fill[targetgreen!22] (0,0) -- (Ai) -- (Aij) -- (Aj) -- cycle;
    \draw[targetgreen!80!black, thick, dashed] (Ai) -- (Aij) -- (Aj);

    \draw[->, very thick, targetgreen!80!black] (0,0) -- (Ai)
      node[below right=2pt, lbl]{$A\hat{i} = (2,1)$};
    \draw[->, very thick, targetgreen!80!black] (0,0) -- (Aj)
      node[above left=2pt, lbl]{$A\hat{j} = (1,2)$};

    \node[lbl, align=center, font=\scriptsize] at (1.5, 1.5) {Area $= \det(A)$\\$= 2\cdot 2 - 1\cdot 1 = 3$};
  \end{scope}

  \node[lbl, align=center] at (3.5, -1.4) {\textbf{(a) Invertible ($\det(A) = 3 \ne 0$)}\\[1pt]Area scales by factor $\det(A)$\\Inverse $A^{-1}$ restores original square};
\end{scope}

% Panel (b): Singular transformation (det = 0)
\begin{scope}[shift={(9.5,0)}]
  \draw[->, gray!60] (-0.3,0) -- (3.4,0) node[right, lbl, black]{\scriptsize $x'$};
  \draw[->, gray!60] (0,-0.3) -- (0,3.2) node[above, lbl, black]{\scriptsize $y'$};

  % Matrix B = [2, 1; 4, 2] -> columns (1.2, 2.4) and (0.6, 1.2)
  \coordinate (Bi) at (1.2, 2.4);
  \coordinate (Bj) at (0.6, 1.2);

  % Collapsed 1D line
  \draw[very thick, trapred!80!black] (0,0) -- (1.5, 3.0)
    node[above right, lbl, trapred!80!black]{line $y = 2x$};

  \draw[->, very thick, trapred] (0,0) -- (Bi)
    node[above left=1pt, lbl, trapred]{$B\hat{i} = (2,4)$};
  \draw[->, very thick, quotientorange] (0,0) -- (Bj)
    node[above left=1pt, lbl, quotientorange]{$B\hat{j} = (1,2)$};

  \node[lbl, align=center, font=\scriptsize, anchor=west] at (1.5, 0.8) {\textbf{Area $= 0$}\\[1pt]Columns dependent};

  \node[lbl, align=center] at (1.7, -1.4) {\textbf{(b) Singular ($\det(B) = 0$)}\\[1pt]2D area squashed to 1D line\\Irreversible: no inverse $B^{-1}$ exists};
\end{scope}
\end{tikzpicture}
\caption{The geometric meaning of the determinant. \textbf{(a)} A linear
transformation maps the unit square to the parallelogram spanned by its columns
$A\hat{i}$ and $A\hat{j}$. The determinant $\det(A) = ad - bc$ is the factor by
which areas scale. \textbf{(b)} When $\det(B) = 0$, the columns are linearly
dependent; 2D area collapses into a 1D line, destroying dimensionality. Because
a flattened line cannot be uniquely unpacked into a plane, singular matrices
cannot be inverted.}
\label{ch39:fig-determinant-area}
\end{figure}

\begin{historybox}
This sliver of linear algebra is CLRS Appendix D. CLRS uses it mainly to support
Chapter 28 (solving systems of linear equations via LU decomposition) and the
matrix view of dynamic programming. If you want the full and beautiful story --
eigenvalues, orthogonality, the spectral theorem -- a dedicated linear-algebra
text is the right next step; we only borrow what the algorithms need.
\end{historybox}

\section{Graphs as formal objects}
\label{ch39:graphs}

So much of this book -- shortest paths, spanning trees, network flow, topological
sorting -- lives on \emph{graphs}. The main chapters use graphs operationally;
here we pin down the definitions precisely, so that when a theorem says ``let
$G = (V, E)$ be a connected, undirected graph,'' every word has met you before.

\subsection{The definition}

\begin{definition}[Graph]
\label{ch39:def-graph}
A \emph{graph} is a pair $G = (V, E)$, where $V$ is a finite set of
\emph{vertices} (also called \emph{nodes}) and $E$ is a set of \emph{edges}. In a
\emph{directed} graph each edge is an ordered pair $(u, v)$ with $u, v \in V$,
drawn as an arrow from $u$ to $v$. In an \emph{undirected} graph each edge is an
unordered pair $\{u, v\}$, drawn as a plain line between $u$ and $v$.
\end{definition}

That is the whole formal content: a set of dots and a set of connections between
them. Everything else -- paths, cycles, connectivity, weights -- is built on top
of this skeleton. We usually write $\lvert V \rvert = n$ for the number of
vertices and $\lvert E \rvert = m$ for the number of edges, and inside asymptotic
notation it is traditional to drop the bars and just write $V$ and $E$, so that
``$O(V + E)$'' means ``$O(\lvert V \rvert + \lvert E \rvert)$.''

\begin{intuitionbox}
An undirected edge is a two-way street: if you can walk from $u$ to $v$, you can
walk back. A directed edge is a one-way street: the arrow $(u,v)$ lets you go from
$u$ to $v$ but says nothing about going back. Almost every modeling decision in
graph algorithms starts with the question, ``are my relationships symmetric?''
Friendship on a typical social network is undirected; ``follows'' on a microblog
is directed.
\end{intuitionbox}

Here is a small directed graph $G = (V, E)$ with vertex set
$V = \{1, 2, 3, 4\}$ and edge set
\[
E = \{(1,2),\ (1,3),\ (2,3),\ (3,4),\ (4,2)\}.
\]

\begin{center}
\begin{tikzpicture}[->,>={Latex[length=2.5mm]},thick,
    node/.style={circle,draw,minimum size=8mm,inner sep=0pt,font=\small}]
  \node[node] (1) at (0,1.6) {1};
  \node[node] (2) at (2,1.6) {2};
  \node[node] (3) at (1,0)   {3};
  \node[node] (4) at (3,0)   {4};
  \draw (1) -- (2);
  \draw (1) -- (3);
  \draw (2) -- (3);
  \draw (3) -- (4);
  \draw (4) -- (2);
\end{tikzpicture}
\end{center}

\subsection{The vocabulary you actually need}

A few terms come up on nearly every page of the graph chapters:

\begin{itemize}
  \item \textbf{Adjacent / neighbor.} Vertex $v$ is adjacent to $u$ if the edge
        between them exists. In a directed graph we are careful about direction.
  \item \textbf{Degree.} In an undirected graph, the \emph{degree} of a vertex is
        the number of edges touching it. In a directed graph we split this into
        \emph{in-degree} (arrows coming in) and \emph{out-degree} (arrows going
        out).
  \item \textbf{Path.} A \emph{path} is a sequence of vertices
        $v_0, v_1, \dots, v_\ell$ where each consecutive pair is joined by an
        edge. Its \emph{length} is the number of edges, $\ell$. A path is
        \emph{simple} if it repeats no vertex.
  \item \textbf{Cycle.} A \emph{cycle} is a path that starts and ends at the same
        vertex (with at least one edge). A graph with no cycles is
        \emph{acyclic}; a directed acyclic graph gets the famous acronym
        \emph{DAG}.
  \item \textbf{Connected.} An undirected graph is \emph{connected} if there is a
        path between every pair of vertices. The maximal connected pieces are its
        \emph{connected components}.
  \item \textbf{Weighted.} A \emph{weight} function $w: E \to \R$ attaches a
        number to each edge -- a distance, a cost, a capacity. Shortest-path and
        minimum-spanning-tree algorithms live on weighted graphs.
\end{itemize}

\begin{formalbox}
\textbf{Handshake lemma.} In any undirected graph, the sum of all vertex degrees
equals twice the number of edges:
\[
\sum_{v \in V} \deg(v) = 2\,\lvert E \rvert.
\]
\emph{Why:} each edge has two endpoints, so it contributes exactly $1$ to the
degree of each of its two vertices -- $2$ in total. Summing over all edges counts
each edge's contribution twice, once per endpoint. A pleasant corollary: the
number of vertices with \emph{odd} degree is always even.
\end{formalbox}

\subsection{Two ways to store a graph}
\label{ch39:graph-rep}

Algorithms do not see your hand-drawn picture; they see a data structure. There
are two standard representations, and choosing between them is a recurring
real-world decision.

\textbf{Adjacency matrix.} Number the vertices $1, \dots, n$ and build the
$n \times n$ matrix $M$ where $M_{ij} = 1$ if the edge $(i,j)$ exists and
$M_{ij} = 0$ otherwise. (For weighted graphs, store the weight instead of $1$.)
For our directed example above:
\[
M =
\begin{array}{c|cccc}
   & 1 & 2 & 3 & 4 \\ \hline
1 & 0 & 1 & 1 & 0 \\
2 & 0 & 0 & 1 & 0 \\
3 & 0 & 0 & 0 & 1 \\
4 & 0 & 1 & 0 & 0
\end{array}
\]
Reading row $1$: vertex $1$ points to $2$ and $3$, matching the edges $(1,2)$ and
$(1,3)$. For an \emph{undirected} graph the matrix is symmetric
($M = M^{\mathsf T}$), because the edge $\{u,v\}$ marks both $M_{uv}$ and
$M_{vu}$.

\textbf{Adjacency list.} For each vertex, store a list of its out-neighbors. Our
example becomes:
\[
1 \to [2, 3], \quad 2 \to [3], \quad 3 \to [4], \quad 4 \to [2].
\]
Pictured as actual linked lists, one row per vertex, the structure looks like
Figure~\ref{ch39:fig-adjlist}: the leftmost column is the array of vertices, and
each vertex points along a chain of its out-neighbors.

\begin{figure}[H]
\begin{center}
\begin{tikzpicture}[
    every node/.style={font=\small},
    vtx/.style={draw,minimum width=7mm,minimum height=7mm,fill=sourceblue!12},
    cell/.style={draw,minimum width=7mm,minimum height=7mm,fill=targetgreen!12},
    nul/.style={font=\scriptsize\itshape},
    >={Latex[length=2mm]}]
  % vertex column
  \node[vtx] (v1) at (0,2.4) {1};
  \node[vtx] (v2) at (0,1.6) {2};
  \node[vtx] (v3) at (0,0.8) {3};
  \node[vtx] (v4) at (0,0.0) {4};
  % chains
  \node[cell] (a1) at (1.4,2.4) {2};
  \node[cell] (a2) at (2.4,2.4) {3};
  \node[cell] (b1) at (1.4,1.6) {3};
  \node[cell] (c1) at (1.4,0.8) {4};
  \node[cell] (d1) at (1.4,0.0) {2};
  % arrows from vertices into chains
  \draw[->] (v1.east) -- (a1.west);
  \draw[->] (a1.east) -- (a2.west);
  \draw[->] (v2.east) -- (b1.west);
  \draw[->] (v3.east) -- (c1.west);
  \draw[->] (v4.east) -- (d1.west);
  % null terminators
  \draw[->] (a2.east) -- ++(0.55,0) node[nul,right] {/};
  \draw[->] (b1.east) -- ++(0.55,0) node[nul,right] {/};
  \draw[->] (c1.east) -- ++(0.55,0) node[nul,right] {/};
  \draw[->] (d1.east) -- ++(0.55,0) node[nul,right] {/};
\end{tikzpicture}
\end{center}
\caption{The adjacency list for the directed example. Each vertex (blue) heads a
chain of its out-neighbors (green); ``/'' marks the end of a chain. Only the five
edges that exist take up space -- contrast the $16$-cell matrix above.}
\label{ch39:fig-adjlist}
\end{figure}

\begin{intuitionbox}
The matrix is a \emph{grid}: instant ``is there an edge from $i$ to $j$?''
lookups, but it always reserves $n^2$ cells whether the graph is dense or nearly
empty. The list is a \emph{phone book}: it stores only the edges that exist, so it
is lean for sparse graphs, but answering ``is $i$ joined to $j$?'' means scanning
$i$'s list. Dense graph, lots of edge queries: matrix. Sparse graph, lots of
neighbor sweeps: list.
\end{intuitionbox}

\begin{complexitybox}
Let $n = \lvert V \rvert$ and $m = \lvert E \rvert$.
\begin{center}
\begin{tabular}{lcc}
\toprule
 & \textbf{Adjacency matrix} & \textbf{Adjacency list} \\
\midrule
Space                       & $\Theta(n^2)$ & $\Theta(n + m)$ \\
Test edge $(u,v)$?          & $\Theta(1)$   & $\Theta(\deg u)$ \\
List neighbors of $u$       & $\Theta(n)$   & $\Theta(\deg u)$ \\
Iterate all edges           & $\Theta(n^2)$ & $\Theta(n + m)$ \\
\bottomrule
\end{tabular}
\end{center}
Most textbook graph algorithms (BFS, DFS, Dijkstra, Prim) assume the adjacency
\emph{list}, because their running times like $O(V + E)$ only emerge when iterating
neighbors is cheap.
\end{complexitybox}

\begin{aestheticbox}
A surprising bonus of the matrix view: the $(i,j)$ entry of $M^k$ (the adjacency
matrix raised to the $k$th power) counts the number of distinct walks of length
exactly $k$ from vertex $i$ to vertex $j$. So a question about \emph{paths} turns
into a question about \emph{matrix powers}, and matrix exponentiation by squaring
answers reachability and counting questions fast. This is the bridge connecting
the first half of this appendix to the second.

Why does it work? The defining sum $(M^2)_{ij} = \sum_{k} M_{ik}\,M_{kj}$ adds a
$1$ for exactly those middle vertices $k$ with both an edge $i \to k$ and an edge
$k \to j$ -- that is, for every two-edge route $i \to k \to j$. For our example,
$(M^2)_{14} = 1$ because the single route $1 \to 3 \to 4$ exists
(Figure~\ref{ch39:fig-walk}), and you can read that very handshake off row $1$ and
column $4$ of $M$.

\begin{figure}[H]
\begin{center}
\begin{tikzpicture}[->,>={Latex[length=2.2mm]},thick,
    node/.style={circle,draw,minimum size=8mm,inner sep=0pt,font=\small}]
  \node[node] (1) at (0,1.6) {1};
  \node[node] (2) at (2,1.6) {2};
  \node[node] (3) at (1,0)   {3};
  \node[node] (4) at (3,0)   {4};
  \draw[sourceblue,line width=1pt] (1) -- (3);
  \draw[sourceblue,line width=1pt] (3) -- (4);
  \draw[gray!55] (1) -- (2);
  \draw[gray!55] (2) -- (3);
  \draw[gray!55] (4) -- (2);
\end{tikzpicture}
\end{center}
\caption{The unique length-$2$ walk from $1$ to $4$, namely $1 \to 3 \to 4$
(highlighted), is exactly what $(M^2)_{14} = 1$ records.}
\label{ch39:fig-walk}
\end{figure}
\end{aestheticbox}

\section{Logarithms, factorials, and identities you will reuse}
\label{ch39:logs-factorials}

The last toolbox is the analytic one: the functions that show up in every running
time. We treat logarithms first because they are the soul of ``divide and
conquer,'' then factorials and a few summation identities.

\subsection{What a logarithm really is}

Start with the friendly definition: $\log_b x$ is the answer to the question
``\emph{to what power must I raise $b$ to get $x$?}'' So
\[
\log_b x = y \quad\Longleftrightarrow\quad b^y = x.
\]
The number $b$ is the \emph{base}. Because $2^3 = 8$, we have $\log_2 8 = 3$.
Because $10^2 = 100$, we have $\log_{10} 100 = 2$. That is the entire idea; the
logarithm is just exponentiation read backwards.

\begin{intuitionbox}
For base $2$ -- the base that dominates this book -- there is an even homelier
reading: $\log_2 x$ is roughly ``how many times can I halve $x$ before I reach
$1$?'' Halve $16$: you get $8, 4, 2, 1$ -- four halvings (Figure~\ref{ch39:fig-loghalve}),
and indeed
$\log_2 16 = 4$. Every time an algorithm cuts its problem in half (binary search,
merge sort, balanced trees), a $\log_2 n$ slips into the running time, and that
``how many halvings'' picture is exactly why.
\end{intuitionbox}

\begin{figure}[H]
\begin{center}
\begin{tikzpicture}[every node/.style={font=\small},
    box/.style={draw,minimum height=7mm,fill=sourceblue!12},
    >={Latex[length=2mm]}]
  \node[box,minimum width=64mm] (b16) at (0,3.0) {$16$};
  \node[box,minimum width=32mm] (b8)  at (-1.6,2.0) {$8$};
  \node[box,minimum width=16mm] (b4)  at (-2.4,1.0) {$4$};
  \node[box,minimum width=8mm]  (b2)  at (-2.8,0.0) {$2$};
  \node[box,minimum width=4mm]  (b1)  at (-3.0,-1.0) {$1$};
  \draw[->] (b16.south west) -- (b8.north west);
  \draw[->] (b8.south west)  -- (b4.north west);
  \draw[->] (b4.south west)  -- (b2.north west);
  \draw[->] (b2.south west)  -- (b1.north west);
  \node[targetgreen,right,font=\footnotesize] at (1.2,2.0) {halving 1};
  \node[targetgreen,right,font=\footnotesize] at (1.2,1.0) {halving 2};
  \node[targetgreen,right,font=\footnotesize] at (1.2,0.0) {halving 3};
  \node[targetgreen,right,font=\footnotesize] at (1.2,-1.0) {halving 4};
  \node[structurepurple,right,font=\footnotesize] at (1.2,3.0)
    {$\log_2 16 = 4$ halvings to reach $1$};
\end{tikzpicture}
\end{center}
\caption{Reading $\log_2 16$ as repeated halving: $16 \to 8 \to 4 \to 2 \to 1$ is
four steps, and indeed $\log_2 16 = 4$. Every ``cut the problem in half'' algorithm
takes about this many steps.}
\label{ch39:fig-loghalve}
\end{figure}

The logarithm grows \emph{achingly} slowly. To make $\log_2 n$ reach $10$ you need
$n = 1024$; to reach $20$ you need $n \approx 10^6$; to reach $30$ you need
$n \approx 10^9$. This is the whole reason an $O(n \log n)$ sort feels almost as
fast as a linear scan, and why an $O(\log n)$ search through a billion items takes
about thirty comparisons.

\subsection{The identities, with quick justifications}

You will use these constantly. Each one is a direct consequence of the
exponent rules, since logs and exponents are inverses.

\begin{formalbox}
For positive reals and a valid base $b > 0, b \ne 1$:
\begin{align*}
\log_b(xy)       &= \log_b x + \log_b y
   &&\text{(a product becomes a sum)}\\
\log_b(x/y)      &= \log_b x - \log_b y
   &&\text{(a quotient becomes a difference)}\\
\log_b(x^k)      &= k\,\log_b x
   &&\text{(an exponent slides out front)}\\
\log_b 1         &= 0, \qquad \log_b b = 1 \\
b^{\log_b x}     &= x, \qquad \log_b(b^k) = k
   &&\text{(inverse functions)}\\
\log_b x         &= \dfrac{\log_a x}{\log_a b}
   &&\text{(change of base)}
\end{align*}
\end{formalbox}

The product rule is the keystone, so let us see why it holds. Write
$x = b^p$ and $y = b^q$, which is just saying $p = \log_b x$ and $q = \log_b y$.
Then $xy = b^p \cdot b^q = b^{p+q}$, and reading that backwards gives
$\log_b(xy) = p + q = \log_b x + \log_b y$. Every other identity falls out the
same way, by translating to exponents and using the exponent laws.

The \emph{change-of-base} formula matters most in an algorithms book.
It says all logarithms are the same up to a constant multiplier, because
$\log_b x = \log_a x / \log_a b$ and $\log_a b$ is a fixed constant once $a$ and
$b$ are chosen. That is why we are cavalier and write ``$O(\log n)$'' without
naming a base: switching from $\log_2$ to $\log_{10}$ to $\ln$ only multiplies by
a constant, and big-$O$ swallows constants whole.

\begin{trapbox}
Constants matter \emph{outside} big-$O$, though. ``Time $= 3 \log_2 n$ seconds''
and ``time $= 3 \ln n$ seconds'' are genuinely different stopwatch readings -- the
second is about $1.44$ times longer, since $\log_2 n = \ln n / \ln 2$ and
$1/\ln 2 \approx 1.4427$. The base is free to ignore only when it lives behind an
$O$, $\Theta$, or $\Omega$.
\end{trapbox}

\subsection{Factorials and how fast they explode}

The \emph{factorial} of a non-negative integer $n$ is
\[
n! = n \cdot (n-1) \cdot (n-2) \cdots 2 \cdot 1,
\qquad\text{with the convention } 0! = 1.
\]
It counts the number of ways to arrange $n$ distinct objects in a row: $n$ choices
for the first slot, $n-1$ for the next, and so on. Watch it run away from you:
\[
\begin{array}{c|cccccccc}
n   & 1 & 2 & 3 & 4 & 5   & 6   & 7    & 10 \\ \hline
n!  & 1 & 2 & 6 & 24 & 120 & 720 & 5040 & 3\,628\,800
\end{array}
\]
By $n = 20$ the factorial has passed $2 \times 10^{18}$. This savage growth is
exactly why brute-force ``try every permutation'' algorithms (the naive
travelling salesman, for instance) are hopeless beyond tiny inputs: $20! $
arrangements would take longer to enumerate than the age of the universe.

\begin{intuitionbox}
Factorial growth is the wall at the end of brute force. Exponential growth
($2^n$) is already a disaster, but factorial growth is \emph{worse than any fixed
exponential}: $n!$ eventually outruns $c^n$ for every constant $c$, because each
new factor $n$ is bigger than any fixed $c$. When you see $n!$ in a running time,
read it as ``this only works for $n$ in the low twenties, ever.''
\end{intuitionbox}

How big is $n!$ exactly? \emph{Stirling's approximation} gives a remarkably tight
handle:
\[
n! \;\approx\; \sqrt{2\pi n}\,\left(\frac{n}{e}\right)^{n}.
\]
Taking logs (which is what we usually want for complexity), it tells us
\[
\log(n!) = \Theta(n \log n).
\]
That single fact, $\log(n!) = \Theta(n\log n)$, is the engine behind the famous
lower bound that \emph{any} comparison-based sorting algorithm needs
$\Omega(n \log n)$ comparisons: there are $n!$ possible orderings to distinguish,
a decision tree distinguishing them needs $\log_2(n!)$ levels, and $\log_2(n!) =
\Theta(n \log n)$.

\subsection{Two summation identities worth memorizing}

When you analyze loops, sums fall out, and two of them are everywhere.

\begin{formalbox}
\textbf{Arithmetic series} (the cost of nested loops where the inner count grows
linearly):
\[
1 + 2 + 3 + \cdots + n = \sum_{k=1}^{n} k = \frac{n(n+1)}{2} = \Theta(n^2).
\]
\textbf{Geometric series} (the cost of recursion that shrinks by a factor each
level), for $r \ne 1$:
\[
1 + r + r^2 + \cdots + r^{n} = \sum_{k=0}^{n} r^k = \frac{r^{n+1} - 1}{r - 1}.
\]
When $\lvert r \rvert < 1$ the infinite version converges:
$\sum_{k=0}^{\infty} r^k = \dfrac{1}{1 - r}$.
\end{formalbox}

The arithmetic series has a one-line proof you can show a child, attributed (maybe
apocryphally) to a young Gauss. Write the sum forwards and backwards and add
column by column:
\[
\begin{array}{ccccccc}
S = & 1 & + & 2 & + \cdots + & n & \\
S = & n & + & (n{-}1) & + \cdots + & 1 & \\ \hline
2S = & (n{+}1) & + & (n{+}1) & + \cdots + & (n{+}1) &
\end{array}
\]
There are $n$ columns, each summing to $n+1$, so $2S = n(n+1)$ and
$S = n(n+1)/2$. This is why a double loop where the inner loop runs $1, 2, 3,
\dots, n$ times costs $\Theta(n^2)$, the recurring fingerprint of selection sort,
insertion sort's worst case, and naive nested scans.

\begin{labbox}
Trace the geometric sum yourself for the case that matters most in this book,
$r = 1/2$:
\[
1 + \tfrac12 + \tfrac14 + \tfrac18 + \cdots = \frac{1}{1 - \tfrac12} = 2.
\]
Now connect it to algorithms. When a recursion does $\Theta(n)$ work at the top,
$\Theta(n/2)$ across the next level, $\Theta(n/4)$ below that, and so on, the
\emph{total} work is $n(1 + \tfrac12 + \tfrac14 + \cdots) < 2n = \Theta(n)$. The
geometric series is precisely why such ``halving'' recursions stay linear instead
of blowing up. You will meet this exact computation when analyzing building a heap
and several divide-and-conquer recurrences.
\end{labbox}

The geometric collapse is easiest to believe as a picture. Stack the per-level
work as bars whose heights halve each time, and watch the running total creep
toward $2n$ without ever reaching it (Figure~\ref{ch39:fig-geometric}).

\begin{figure}[H]
\begin{center}
\begin{tikzpicture}[scale=0.9, every node/.style={font=\footnotesize}]
  \def\h{1.6}
  % baseline
  \draw[->,thick] (-0.3,0) -- (7.0,0) node[right] {level};
  % bars: widths 1, heights n, n/2, n/4, ...
  \fill[sourceblue!70] (0,0) rectangle (1,\h);
  \fill[sourceblue!60] (1.2,0) rectangle (2.2,\h/2);
  \fill[sourceblue!50] (2.4,0) rectangle (3.4,\h/4);
  \fill[sourceblue!40] (3.6,0) rectangle (4.6,\h/8);
  \fill[sourceblue!30] (4.8,0) rectangle (5.8,\h/16);
  \node at (0.5,\h+0.25) {$n$};
  \node at (1.7,\h/2+0.25) {$\tfrac{n}{2}$};
  \node at (2.9,\h/4+0.25) {$\tfrac{n}{4}$};
  \node at (4.1,\h/8+0.28) {$\tfrac{n}{8}$};
  \node at (5.3,\h/8) {$\cdots$};
  % running total marker
  \draw[densely dashed,targetgreen,thick] (-0.3,\h*1.92) -- (6.4,\h*1.92);
  \node[targetgreen,above left] at (6.4,\h*1.92) {total $\to 2n$};
\end{tikzpicture}
\end{center}
\caption{Per-level work in a ``halving'' recursion. Each bar is half the height of
the one before, so the heights sum to $n(1+\tfrac12+\tfrac14+\cdots)=2n$. The dashed
line is the ceiling the running total approaches but never crosses.}
\label{ch39:fig-geometric}
\end{figure}

\begin{historybox}
This section maps to CLRS Appendix A (summations) and the logarithm/factorial
facts scattered through Appendix C and Chapter 3 (asymptotic notation).
Stirling's approximation and the $\log(n!) = \Theta(n\log n)$ consequence anchor
the comparison-sort lower bound in CLRS Chapter 8, and the geometric-series
collapse is the workhorse behind the master theorem in CLRS Chapter 4.
\end{historybox}

\section{Putting the toolbox together}
\label{ch39:wrap}

Step back and notice how interlocked these ``separate'' tools really are. The
adjacency \emph{matrix} let us turn a question about graph \emph{paths} into a
question about matrix \emph{powers}. Computing those powers fast uses
\emph{exponentiation by squaring}, whose efficiency rests on a
\emph{logarithm} -- you square $\log_2 k$ times to reach the $k$th power.
Counting the arrangements a sorting algorithm must distinguish brought in
\emph{factorials}, and pinning down how big a factorial is brought back the
\emph{logarithm} through Stirling. The discrete and the linear-algebraic are not
two worlds; they are two views of the same machinery for reasoning about
structure and growth.

\begin{focusbox}
The four things to carry out of this appendix:
\begin{enumerate}
  \item A matrix is a spreadsheet; multiplication is row-meets-column handshakes,
        costs $\Theta(n^3)$ naively, and is \emph{not} commutative.
  \item A graph is $(V, E)$, stored as a matrix ($\Theta(n^2)$, fast edge tests)
        or a list ($\Theta(n+m)$, fast neighbor sweeps).
  \item $\log_b x$ answers ``$b$ to what power is $x$?''; it grows glacially,
        turns products into sums, and is base-independent under big-$O$.
  \item $n!$ explodes faster than any exponential, and $\log(n!) =
        \Theta(n \log n)$ is the seed of the comparison-sort lower bound.
\end{enumerate}
\end{focusbox}

\subsection{Controversies and false starts: non-commutativity and numerical collapse}

The abstraction of matrices required overcoming long-standing mathematical
prejudices and dangerous numerical pitfalls:
\begin{itemize}
  \item \textbf{The shock of non-commutativity.} When Arthur Cayley published his
        theory of matrices in 1858, 19th-century mathematicians were deeply
        disquieted by the fact that matrix multiplication is non-commutative:
        $AB \ne BA$ in general. For thousands of years, arithmetic had rested on
        the commutative truth that $a \times b = b \times a$. An operation where
        order matters seemed unnatural and hazardous to many traditionalists,
        delaying the widespread adoption of matrix algebra until its essential role
        in quantum mechanics (Heisenberg's matrix mechanics in 1925) proved its
        necessity.
  \item \textbf{The Leibniz expansion trap for determinants.} Early textbooks taught
        students to evaluate determinants via Leibniz's permutation formula
        $\sum_\sigma \operatorname{sgn}(\sigma) \prod A_{i,\sigma(i)}$ or Laplace's
        cofactor expansion. Both require $n!$ operations. For an ordinary
        $30 \times 30$ matrix, $30! \approx 2.65 \times 10^{32}$ operations---taking
        longer than the age of the universe on the fastest supercomputer!
        Gaussian elimination evaluates the exact same determinant in $\Theta(n^3)$
        steps by converting the matrix to upper-triangular form and multiplying the
        diagonal elements.
  \item \textbf{The disaster of naive Gaussian elimination.} In textbook mathematics,
        Gaussian elimination always works if the matrix is non-singular. On a
        computer using floating-point numbers, however, naive elimination without
        \emph{pivoting} is an engineering disaster. If a diagonal pivot element
        $A_{ii}$ happens to be very small, dividing by it magnifies rounding
        errors exponentially, causing subsequent row operations to completely
        corrupt the solution vector. The invention of \emph{partial pivoting}
        (swapping rows to place the largest available entry on the diagonal) was
        required to turn elimination into a numerically stable algorithm.
\end{itemize}

\subsection{What this unlocked: the universal computational substrate}

Matrix algebra provided the unified language across science, engineering, and
computer science:
\begin{itemize}
  \item \textbf{High-dimensional geometric transformations.} In 3D computer graphics,
        virtual reality, and robotics, $4 \times 4$ homogeneous transformation
        matrices represent rotations, translations, perspective projections, and
        scaling. Hardware GPUs execute billions of these matrix-vector transforms
        every second.
  \item \textbf{Graph theory as algebra.} Encoding graph adjacency in matrices
        turns combinatorial questions into algebraic ones. Powers of the
        adjacency matrix $A^k$ count paths, the graph Laplacian $L = D - A$
        measures cuts and electrical flows, and the principal eigenvector of the
        Google matrix yields the PageRank of every page on the Web.
  \item \textbf{The foundation of artificial intelligence.} Every modern deep neural
        network is an alternating composition of matrix multiplications (weight
        tensors) and point-wise non-linear activations. Hardware architectures
        like TPUs and Tensor Cores are specialized matrix engines built
        specifically to multiply Cayley's grids at silicon speed.
\end{itemize}

\subsection{Honest downsides: the cubic wall and sparsity demands}

Despite their ubiquity, dense matrix operations face severe algorithmic and physical
limitations:
\begin{itemize}
  \item \textbf{The cubic complexity bottleneck.} Naive matrix multiplication costs
        $\Theta(n^3)$ operations. While sub-cubic algorithms exist---from Volker
        Strassen's $O(n^{2.807})$ divide-and-conquer algorithm (1969) to modern
        laser methods approaching $O(n^{2.37})$---their immense constant factors
        and numerical instability mean practical libraries (BLAS/LAPACK) still
        rely heavily on blocked $O(n^3)$ algorithms optimized for hardware caches.
  \item \textbf{Memory quadratic explosion and sparsity.} Storing an $n \times n$
        dense matrix requires $\Theta(n^2)$ memory. For large real-world graphs,
        such as the social network graph with $n = 3 \times 10^9$ users, a dense
        matrix would require tens of exabytes of RAM! Systems must use sparse
        representations (like Compressed Sparse Row, CSR) that store only the
        non-zero entries, sacrificing fast $O(1)$ coordinate access.
  \item \textbf{Ill-conditioning and numerical instability.} If a matrix has a large
        condition number $\kappa(A) = \|A\| \|A^{-1}\|$, solving the linear
        system $Ax = b$ is fundamentally ill-conditioned: tiny roundoff errors or
        measurement noise in $b$ produce massive, wildly inaccurate errors in the
        computed solution $x$.
\end{itemize}

\subsection{Forward roadmap: the journey through the curriculum}

With this appendix, you have completed the foundational journey of
\emph{Algorithms From the Ground Up}. The tools forged across these thirty-nine
chapters open directly into the advanced volumes of this curriculum:
\begin{itemize}
  \item \textbf{Linear Algebra From the Ground Up (Book~3):} Explores vector
        spaces, orthogonal projections, eigenvalues, eigenvectors, Singular
        Value Decomposition (SVD), and principal component analysis (PCA) from
        first principles.
  \item \textbf{Graph Algorithms From the Ground Up (Book~2):} Deepens the study
        of network flows, bipartite matchings, planar graphs, tree decompositions,
        and spectral graph theory using the graph Laplacian.
  \item \textbf{Signal Processing From the Ground Up (Book~4):} Applies linear
        transformations and the Discrete Fourier Transform matrix to sound,
        images, and continuous time series.
  \item \textbf{Machine Learning From the Ground Up:} Builds high-dimensional
        optimization, backpropagation, and transformer architectures on the
        matrix foundations assembled here.
\end{itemize}

\section*{Exercises}
\addcontentsline{toc}{section}{Exercises}

\begin{exercise}[Warmup: shapes]
For each pair, decide whether the product $AB$ exists and, if so, give the shape
of the result: (a) $A$ is $3 \times 4$, $B$ is $4 \times 2$; (b) $A$ is
$2 \times 5$, $B$ is $2 \times 5$; (c) $A$ is $1 \times n$, $B$ is $n \times 1$;
(d) $A$ is $n \times 1$, $B$ is $1 \times n$. For (c) and (d), what kind of object
is the result?
\end{exercise}

\begin{exercise}[Multiply by hand]
Let $A = \begin{psmallmatrix} 2 & -1 \\ 0 & 3 \end{psmallmatrix}$ and
$B = \begin{psmallmatrix} 1 & 4 \\ 2 & 1 \end{psmallmatrix}$. Compute both $AB$
and $BA$ entry by entry and confirm they differ. Then compute $A^{\mathsf T}$ and
verify $(AB)^{\mathsf T} = B^{\mathsf T} A^{\mathsf T}$.
\end{exercise}

\begin{exercise}[Identity and inverse]
Verify that $I_2 A = A$ for the matrix $A$ of the previous exercise. Then compute
$\det(B)$ and, if it is nonzero, write down $B^{-1}$ using the $2 \times 2$
formula and check that $B B^{-1} = I_2$.
\end{exercise}

\begin{exercise}[Degrees and the handshake]
Draw any undirected graph on $5$ vertices with exactly $6$ edges. Write down the
degree of each vertex, sum them, and confirm the sum equals $2 \times 6 = 12$.
Then explain in one sentence why a graph in which every one of the $5$ vertices
has odd degree is impossible.
\end{exercise}

\begin{exercise}[Two representations]
For the directed graph with $V = \{a,b,c,d\}$ and edge set
\[
E = \{(a,b),\ (b,c),\ (c,a),\ (c,d)\},
\]
write both the $4 \times 4$ adjacency matrix and the adjacency lists. Which
representation uses less space here, and why?
\end{exercise}

\begin{exercise}[Walks via matrix powers]
Using the adjacency matrix $M$ of the graph in the previous exercise, compute
$M^2$ and interpret the $(a,c)$ entry: how many length-$2$ walks go from $a$ to
$c$? Trace such a walk in the picture to confirm.
\end{exercise}

\begin{exercise}[Log gymnastics]
Without a calculator, evaluate: (a) $\log_2 64$; (b) $\log_2 (1/8)$;
(c) $\log_2(32 \cdot 4)$ using the product rule; (d) $\log_{10} 10^7$;
(e) $\log_2 1000$ to the nearest integer, given $\log_{10} 2 \approx 0.301$ and
the change-of-base formula.
\end{exercise}

\begin{exercise}[Factorial intuition]
(a) Compute $8!$ by hand. (b) Using $0! = 1$, explain why $n! = n \cdot (n-1)!$ is
a sensible recursive definition all the way down to $n = 1$. (c) Roughly how many
permutations does $13!$ give, and why does that doom a brute-force ordering of
just $13$ tasks?
\end{exercise}

\begin{exercise}[Sum it up]
Use the arithmetic-series formula to find $1 + 2 + \cdots + 100$. Then a loop runs
its body once on the first iteration, twice on the second, $\dots$, $n$ times on
the $n$th. Give a closed form for the total number of body executions and its
$\Theta$ class.
\end{exercise}

\begin{exercise}[Geometric collapse]
Evaluate the finite sum $1 + 3 + 9 + 27 + 81$ with the geometric formula
($r = 3$), then check by direct addition. Separately, argue that
$1 + \tfrac13 + \tfrac19 + \cdots$ converges and find its value.
\end{exercise}

\begin{exercise}[Challenge: cost of a matrix chain]
You must compute the product $A_1 A_2 A_3$ where $A_1$ is $10 \times 100$, $A_2$
is $100 \times 5$, and $A_3$ is $5 \times 50$. Using the rule that multiplying a
$p \times q$ matrix by a $q \times r$ matrix costs $p \cdot q \cdot r$ scalar
multiplications, compute the total cost of $(A_1 A_2) A_3$ versus
$A_1 (A_2 A_3)$. Which parenthesization is cheaper, and by how much? (This is the
seed of the matrix-chain dynamic-programming problem.)
\end{exercise}

\begin{exercise}[Challenge: Stirling and sorting]
Stirling gives $\log_2(n!) = \Theta(n \log n)$. A comparison sort can be modeled
as a binary decision tree whose leaves are the $n!$ possible orderings. Explain
why such a tree must have height at least $\log_2(n!)$, and conclude that every
comparison-based sort makes $\Omega(n \log n)$ comparisons in the worst case.
\end{exercise}
